% \section{Setting and Data}

% \paragraph{Facts and ontology.}
% We use drug--disease edges from PrimeKG \citep{chandak2023primekg}, whose
% indication and contraindication relations come from DrugCentral. Diseases are
% identified by MONDO \citep{vasilevsky2022mondo}, which we rebuild as a directed
% child-to-parent table of 17,801 edges over 11,788 diseases; 3,182 parent terms
% have two or more children.

% \paragraph{Text.}
% For each drug we use its English Wikipedia article and a set of PubMed
% abstracts retrieved per drug with a separate query budget for each relation.
% Wikipedia markup is cleaned before use: templates, tables, reference tags and
% file links are removed, and links are reduced to their visible text.

% \paragraph{How a page names a fact.}
% For each edge we classify the drug's text into one of three groups.
% \emph{Named}: the text contains the disease's exact name or an alternative
% name taken from Wikipedia redirects to the disease's article, which include
% synonyms, lay terms and abbreviations. \emph{Family-only}: the text contains
% none of these but names one of the disease's MONDO ancestors within two
% levels. \emph{Silent}: the text names neither. To separate
% real mentions from coincidence, each edge is paired with three diseases the
% drug has no relation to, and the same matcher is run on them.

% \subsection{How Much Does Text Reach?}

% We first measure, on 1,800 PrimeKG edges (600 per relation for indication,
% contraindication and off-label use), how often the drug's Wikipedia article
% names the disease and how often a model knows the fact closed book. The model
% picks the true disease from five options by mean token log-probability.

% \begin{table}[t]
% \centering
% \small
% \begin{tabular}{lcc}
% \toprule
%  & Model knows it & Model does not \\
% \midrule
% Text names it & 5.2 & 6.2 \\
% Text does not & 22.2 & 66.4 \\
% \bottomrule
% \end{tabular}
% \caption{Share of 1,800 PrimeKG facts (\%) by whether Wikipedia names the
% disease and whether Qwen2.5 0.5B picks it from five options closed book.}
% \label{tab:boxes}
% \end{table}

% Table~\ref{tab:boxes} crosses the two; here ``names it'' means the exact or an
% alternative name, so facts named only by their family count as not named. For
% two thirds of the facts, the text does not name the disease and the model does
% not know the answer. The ``knows it anyway'' cell is mostly chance: accuracy on
% facts the text does not name is 25.1\% against a 20\% chance rate, about 6\%
% above chance after correction. A second corpus helps little: on a separate
% sample measured in both corpora, the share not named is 84.9\% in Wikipedia,
% 81.1\% in PubMed abstracts and 77.5\% in both.

% \paragraph{A family name alone does not help the model.}
% If a family name carried useful signal, facts whose text names only the family
% would be easier for the model than facts whose text is silent. Testing this
% inside each relation and pooling with relation as a stratum on Qwen2.5-7B,
% family-only facts differ sharply from named facts (Mantel--Haenszel odds ratio
% 0.27, 95\% CI 0.17--0.43, $p = 9\times10^{-9}$) and not measurably from
% silent facts (odds ratio 1.10, 95\% CI 0.80--1.53, $p = 0.61$). The cleanest
% case is indications at 7B: 44.68\% correct for family-only facts against
% 43.95\% for silent facts. For contraindications the figures are 18.75\% and
% 23.35\%. An exact name raises the odds that the model knows a fact about
% 3.7-fold; a family name gives no measurable help on its own.

% This motivates the rest of the paper. The family name is the most common form
% of partial evidence in this data, and the model cannot use it by itself. The
% question is what happens when such a page is placed in front of the model at
% answer time, and whether a symbolic step can make use of it.



% \section{Data and Problem Setup}

% \paragraph{Knowledge graphs and ontologies.}
% Our primary benchmark uses drug-disease relations from PrimeKG \citep{chandak2023primekg}, which draws clinical indication and contraindication facts from DrugCentral. Each disease is linked to the MONDO Disease Ontology \citep{vasilevsky2022mondo}. We reconstruct MONDO as a directed child-to-parent hierarchy, with \num{17801} edges across \num{11788} diseases, in which \num{3182} parent categories have two or more child diseases. We exclude parents with more than 400 descendants (e.g. \emph{disease} or \emph{disorder}), since naming them gives no usable hint. To test whether our findings generalize beyond medicine, we also evaluate species-interaction facts from Global Biotic Interactions (GloBI; \citealt{poelen2014globi}). In GloBI, we examine four ecological relations (\emph{eats}, \emph{preys on}, \emph{parasite of}, and \emph{pollinates}) and use the Tree-of-Life biological taxonomy (\emph{species}, \emph{genus}, \emph{family}) as the hierarchy. We keep only species-level names on both sides of an edge, which leaves \num{457386} species-to-species edges. Table~\ref{tab:data} summarizes both domains. The two hierarchies have the same shape: the median family has three members, so a family name narrows the answer to a short list but does not settle it.

% \begin{table}[t]
\centering
\small
\setlength{\tabcolsep}{3.5pt}
\begin{tabular}{@{}lrr@{}}
\toprule
 & PrimeKG & GloBI \\
\midrule
Source records & 8,100,498 & 7,822,670 \\
\midrule
\multirow{4}{*}{Edges} & indication 9,388 & eats 235,960 \\
 & contraind.\ 30,675 & preys on 87,678 \\
 & off-label 2,568 & parasite of 117,894 \\
 & & pollinates 15,854 \\
\midrule
Hierarchy & MONDO & GloBI taxonomy \\
\quad child--parent links & 17,801 & 77,852 \\
\quad families, $\ge$2 members & 3,182 & 8,370 \\
\quad median / max size & 3 / 150 & 3 / 315 \\
\midrule
Text & Wiki.\ + PubMed & Wikipedia \\
\quad sources with text & 574 drugs & 5,425 species \\
\quad facts measured & 3,444 & 23,694 \\
Questions & 856 & 900 \\
\bottomrule
\end{tabular}
\caption{Data for both domains. PrimeKG links drugs to diseases, and MONDO groups diseases into families. GloBI links species to species, and its taxonomy groups species into genera. A \emph{family} is a parent with two or more children. \emph{Facts measured} are edges we check for a mention of the answer.}
\label{tab:data}
\end{table}


% \paragraph{Reference text corpora.}
% For each drug in PrimeKG, we collect its English Wikipedia article and a set of PubMed abstracts retrieved using separate query budgets for indications and contraindications. For each organism in GloBI, we use the consumer species' English Wikipedia article, skipping stubs shorter than 3,000 characters. All Wikipedia articles are stripped of markup, templates, infoboxes, tables, reference tags, and file links, leaving only the plain article text.

% \paragraph{Evidence granularity.}
% For any given fact (e.g. a drug treating a specific disease), we classify the retrieved text into one of three granularity levels:
% \begin{itemize}
% \item \textbf{Named (exact-match):} The passage contains the exact name of the target disease (or target species), or a known synonym, or abbreviation from Wikipedia redirect titles.
% \item \textbf{Family-only (coarse):} The passage does not contain the exact answer or synonyms, but mentions a broader parent category within two levels of the answer in the hierarchy.
% \item \textbf{Silent (unmentioned):} The passage mentions neither the exact answer nor any parent category within two levels.
% \end{itemize}

% To verify that matches reflect genuine mentions rather than chance word overlap, we pair each fact with three unrelated diseases (or species) from the same relation and run the same text matcher as a baseline control. On a separate sample of 1,800 PrimeKG facts measured in Wikipedia, the unrelated diseases are named in 0.9-1.4\% of articles and named or matched by family in 1.8-3.3\%, against 6.7-19.3\% and 9.3-34.7\% for the true disease across the three relations. Matches therefore track the fact, not the size of the name set.

% \paragraph{How often does text have the exact answer?}

% \begin{table}[t]
\centering
\small
\setlength{\tabcolsep}{4pt}
\begin{tabular}{@{}lrrrrr@{}}
\toprule
 & & \multicolumn{2}{c}{Text names it} & \multicolumn{2}{c}{Text does not} \\
\cmidrule(lr){3-4}\cmidrule(l){5-6}
Relation & $n$ & knows & not & knows & not \\
\midrule
Indication & 600 & 9.7 & 7.2 & 21.8 & 61.3 \\
Contraind. & 600 & 1.5 & 5.2 & 20.7 & 72.7 \\
Off-label & 600 & 4.3 & 6.2 & 24.2 & 65.3 \\
\midrule
All & 1,800 & 5.2 & 6.2 & 22.2 & \textbf{66.4} \\
\bottomrule
\end{tabular}
\caption{Share of PrimeKG facts (\%) by whether the drug's Wikipedia article
names the disease and whether Qwen2.5 0.5B picks it from five options closed
book (``knows'').}
\label{tab:boxes}
\end{table}


% Before testing RAG, we test how often reference texts contain the fine-grained answer and how often an LLM knows the fact without any retrieval (closed book). We sample 1,800 PrimeKG facts (600 each for indication, contraindication, and off-label use) and test whether Qwen2.5 0.5B can pick the true disease from five options based on mean token log-probability. ~\cref{tab:boxes} cross-tabulates closed-book model accuracy against whether the drug's Wikipedia page names the exact disease. Here, coarse (\emph{family-only}) mentions are grouped under \emph{Text does not name it}. For two-thirds of the facts (66.4\%), Wikipedia does not state the exact disease and the model does not know the answer. The share is highest for contraindications (72.7\%), where Wikipedia names the disease for only 6.7\% of facts. Even when the model selects the right option without text (22.2\%), much of that success is random guessing: closed-book accuracy on unnamed facts is 25.1\% compared to a 20.0\% five-choice guessing rate, which corresponds to only about 5\% above-chance knowledge.

% \begin{table}[t]
\centering
\small
\setlength{\tabcolsep}{4pt}
\begin{tabular}{@{}llrrrr@{}}
\toprule
Relation & Text & $n$ & Named & Family & Silent \\
\midrule
\multirow{3}{*}{Indication} & Wikipedia & 1,722 & 20.3 & 15.5 & 64.2 \\
 & PubMed & 1,722 & 26.3 & 18.6 & 55.1 \\
 & both & 1,722 & 29.0 & 19.7 & 51.3 \\
\midrule
\multirow{3}{*}{Contraind.} & Wikipedia & 1,722 & 9.8 & 4.1 & 86.1 \\
 & PubMed & 1,722 & 11.4 & 5.4 & 83.2 \\
 & both & 1,722 & 15.9 & 6.9 & 77.2 \\
\midrule
PrimeKG & both & 3,444 & 22.5 & 13.3 & 64.2 \\
\midrule
Eats & Wikipedia & 8,000 & 4.0 & 6.6 & 89.3 \\
Preys on & Wikipedia & 8,000 & 2.2 & 4.4 & 93.4 \\
Parasite of & Wikipedia & 2,389 & 12.5 & 15.5 & 72.0 \\
Pollinates & Wikipedia & 5,305 & 2.2 & 8.0 & 89.7 \\
\midrule
GloBI & Wikipedia & 23,694 & 3.9 & 7.1 & 89.1 \\
\bottomrule
\end{tabular}
\caption{Share of knowledge-graph facts whose text names the answer, names only its family, or names neither. PrimeKG uses 574 drugs with three facts of each relation.}
\label{tab:coverage}
\end{table}


% Adding a second biomedical corpus (\emph{PubMed}) barely changes this. Table~\ref{tab:coverage} measures a drug-matched sample of 574 drugs, each with three indications and three contraindications. The exact disease is missing in 84.9\% of Wikipedia articles, 81.1\% of PubMed abstract sets, and 77.5\% of both corpora combined. Indications are named about twice as often as contraindications in every corpus (29.0\% vs.\ 15.9\% in both combined). Family-only text covers 13.3\% of PrimeKG facts, more than half as many as are named (22.5\%). In GloBI, exact species mentions are even rarer (3.9\%), while coarse genus or family mentions account for nearly twice as many facts (7.1\%). The pattern holds in every GloBI relation, and \emph{parasite of} has the most family-only text (15.5\%).

% \paragraph{Coarse mentions do not reflect stronger closed-book knowledge.}
% One might wonder whether facts described at the family level in text are simply easier facts that the model already knows. If a coarse mention reflected easier underlying facts, then models tested closed book should perform better on family-only facts than on silent facts. We test this on Qwen2.5-7B across all relations using Mantel--Haenszel stratified odds ratios (Table~\ref{tab:closedbook}). Closed-book accuracy on family-only facts is much lower than on named facts (odds ratio $0.27$, 95\% CI $0.17$--$0.43$, $p = 9\times10^{-9}$) and is statistically indistinguishable from silent facts (odds ratio $1.10$, 95\% CI $0.80$--$1.53$, $p = 0.61$). For drug indications, Qwen2.5-7B answers 44.7\% of family-only facts correctly closed book versus 44.0\% of silent facts; for contraindications, the rates are 18.8\% versus 23.4\%. While an exact mention in Wikipedia corresponds to a 3.7-fold increase in the odds that the model knows the fact closed book, a coarse family mention provides no measurable advantage on its own.

% \begin{table}[t]
\centering
\small
\setlength{\tabcolsep}{4pt}
\begin{tabular}{@{}lrrr@{}}
\toprule
Relation & Named & Family only & Silent \\
\midrule
Indication & 84.2 {\scriptsize(101)} & 44.7 {\scriptsize(94)} & 44.0 {\scriptsize(405)} \\
Contraind. & 32.5 {\scriptsize(40)} & 18.8 {\scriptsize(16)} & 23.4 {\scriptsize(544)} \\
Off-label & 57.1 {\scriptsize(63)} & 37.5 {\scriptsize(72)} & 31.6 {\scriptsize(465)} \\
\midrule
All & 65.7 {\scriptsize(204)} & 39.6 {\scriptsize(182)} & 32.0 {\scriptsize(1,414)} \\
\bottomrule
\end{tabular}
\caption{Closed-book accuracy of Qwen2.5-7B by what the drug's Wikipedia page says about the answer. Chance is 20\%.}
\label{tab:closedbook}
\end{table}


% This observation motivates the rest of our study. Coarse category mentions are the most common form of partial evidence in reference text, yet they do not correspond to facts the model already knows. We therefore ask what happens when a coarse, family-only passage is shown to the model at inference time, and whether a symbolic ontology can help recover the exact answer.



% \section{Data and Problem Setup}

% We use drug-disease facts from PrimeKG, linking each disease to the MONDO ontology. We rebuild MONDO as a hierarchy of 17,801 edges across 11,788 diseases, excluding parents with over 400 descendants. To test generalizability, we also evaluate species interactions from GloBI, keeping 457,386 species-to-species edges across four ecological relations. For each entity, we collect its English Wikipedia article (and PubMed abstracts for drugs) and strip all markup.

% \begin{table}[t]
\centering
\small
\setlength{\tabcolsep}{3.5pt}
\begin{tabular}{@{}lrr@{}}
\toprule
 & PrimeKG & GloBI \\
\midrule
Source records & 8,100,498 & 7,822,670 \\
\midrule
\multirow{4}{*}{Edges} & indication 9,388 & eats 235,960 \\
 & contraind.\ 30,675 & preys on 87,678 \\
 & off-label 2,568 & parasite of 117,894 \\
 & & pollinates 15,854 \\
\midrule
Hierarchy & MONDO & GloBI taxonomy \\
\quad child--parent links & 17,801 & 77,852 \\
\quad families, $\ge$2 members & 3,182 & 8,370 \\
\quad median / max size & 3 / 150 & 3 / 315 \\
\midrule
Text & Wiki.\ + PubMed & Wikipedia \\
\quad sources with text & 574 drugs & 5,425 species \\
\quad facts measured & 3,444 & 23,694 \\
Questions & 856 & 900 \\
\bottomrule
\end{tabular}
\caption{Data for both domains. PrimeKG links drugs to diseases, and MONDO groups diseases into families. GloBI links species to species, and its taxonomy groups species into genera. A \emph{family} is a parent with two or more children. \emph{Facts measured} are edges we check for a mention of the answer.}
\label{tab:data}
\end{table}


% We classify retrieved text into three granularity levels:
% \begin{itemize}
% \item \textbf{Named:} The passage contains the exact name or a known redirect synonym.
% \item \textbf{Family-only:} The passage lacks the exact name but mentions a parent category within two levels above in the hierarchy.
% \item \textbf{Silent:} The passage mentions neither.
% \end{itemize}

% Control tests with unrelated disease pairs confirm that text matches track the actual fact rather than name-set size (false matches occur in only 0.9--1.4\% of articles).

% We test how often text states the exact answer and how well Qwen2.5 0.5B knows facts closed book using 1,800 PrimeKG facts. As shown in Table~\ref{tab:boxes} and Table~\ref{tab:coverage}, Wikipedia completely misses the exact disease for 84.9\% of facts (77.5\% across both Wikipedia and PubMed). For two-thirds of facts (66.4\%), the text is missing and the model does not know the answer. 

% \begin{table}[t]
\centering
\small
\setlength{\tabcolsep}{4pt}
\begin{tabular}{@{}lrrrrr@{}}
\toprule
 & & \multicolumn{2}{c}{Text names it} & \multicolumn{2}{c}{Text does not} \\
\cmidrule(lr){3-4}\cmidrule(l){5-6}
Relation & $n$ & knows & not & knows & not \\
\midrule
Indication & 600 & 9.7 & 7.2 & 21.8 & 61.3 \\
Contraind. & 600 & 1.5 & 5.2 & 20.7 & 72.7 \\
Off-label & 600 & 4.3 & 6.2 & 24.2 & 65.3 \\
\midrule
All & 1,800 & 5.2 & 6.2 & 22.2 & \textbf{66.4} \\
\bottomrule
\end{tabular}
\caption{Share of PrimeKG facts (\%) by whether the drug's Wikipedia article
names the disease and whether Qwen2.5 0.5B picks it from five options closed
book (``knows'').}
\label{tab:boxes}
\end{table}

% \begin{table}[t]
\centering
\small
\setlength{\tabcolsep}{4pt}
\begin{tabular}{@{}llrrrr@{}}
\toprule
Relation & Text & $n$ & Named & Family & Silent \\
\midrule
\multirow{3}{*}{Indication} & Wikipedia & 1,722 & 20.3 & 15.5 & 64.2 \\
 & PubMed & 1,722 & 26.3 & 18.6 & 55.1 \\
 & both & 1,722 & 29.0 & 19.7 & 51.3 \\
\midrule
\multirow{3}{*}{Contraind.} & Wikipedia & 1,722 & 9.8 & 4.1 & 86.1 \\
 & PubMed & 1,722 & 11.4 & 5.4 & 83.2 \\
 & both & 1,722 & 15.9 & 6.9 & 77.2 \\
\midrule
PrimeKG & both & 3,444 & 22.5 & 13.3 & 64.2 \\
\midrule
Eats & Wikipedia & 8,000 & 4.0 & 6.6 & 89.3 \\
Preys on & Wikipedia & 8,000 & 2.2 & 4.4 & 93.4 \\
Parasite of & Wikipedia & 2,389 & 12.5 & 15.5 & 72.0 \\
Pollinates & Wikipedia & 5,305 & 2.2 & 8.0 & 89.7 \\
\midrule
GloBI & Wikipedia & 23,694 & 3.9 & 7.1 & 89.1 \\
\bottomrule
\end{tabular}
\caption{Share of knowledge-graph facts whose text names the answer, names only its family, or names neither. PrimeKG uses 574 drugs with three facts of each relation.}
\label{tab:coverage}
\end{table}


% Crucially, family-only mentions do not reflect stronger closed-book knowledge. Tested on Qwen2.5-7B, closed-book accuracy on family-only facts is much lower than on named facts (odds ratio $0.27$, $p = 9\times10^{-9}$) and indistinguishable from silent facts (odds ratio $1.10$, $p = 0.61$). While an exact name in Wikipedia increases the odds that a model knows a fact 3.7-fold, a family name provides no measurable advantage on its own. 

% \begin{table}[t]
\centering
\small
\setlength{\tabcolsep}{4pt}
\begin{tabular}{@{}lrrr@{}}
\toprule
Relation & Named & Family only & Silent \\
\midrule
Indication & 84.2 {\scriptsize(101)} & 44.7 {\scriptsize(94)} & 44.0 {\scriptsize(405)} \\
Contraind. & 32.5 {\scriptsize(40)} & 18.8 {\scriptsize(16)} & 23.4 {\scriptsize(544)} \\
Off-label & 57.1 {\scriptsize(63)} & 37.5 {\scriptsize(72)} & 31.6 {\scriptsize(465)} \\
\midrule
All & 65.7 {\scriptsize(204)} & 39.6 {\scriptsize(182)} & 32.0 {\scriptsize(1,414)} \\
\bottomrule
\end{tabular}
\caption{Closed-book accuracy of Qwen2.5-7B by what the drug's Wikipedia page says about the answer. Chance is 20\%.}
\label{tab:closedbook}
\end{table}


% Coarse mentions are common partial evidence, which prompts us to test how models handle them at inference time and whether a symbolic ontology can recover the exact answer.

\section{Data and Problem Setup}

\paragraph{Knowledge graphs and ontologies.}
We use drug-disease relations from PrimeKG, linked to the MONDO Disease Ontology. We reconstruct MONDO into 17,801 edges across 11,788 diseases (excluding parents with $>$\num{400} descendants). To test generalizability beyond medicine, we evaluate species interactions from GloBI, using 457,386 species-to-species edges across four relations (\emph{eats}, \emph{preys on}, \emph{parasite of}, \emph{pollinates}) with GloBI's taxonomy (species, genus, family) as the hierarchy.~\cref{tab:data} summarizes both domains.

\begin{table}[t]
\centering
\small
\setlength{\tabcolsep}{3.5pt}
\begin{tabular}{@{}lrr@{}}
\toprule
 & PrimeKG & GloBI \\
\midrule
Source records & 8,100,498 & 7,822,670 \\
\midrule
\multirow{4}{*}{Edges} & indication 9,388 & eats 235,960 \\
 & contraind.\ 30,675 & preys on 87,678 \\
 & off-label 2,568 & parasite of 117,894 \\
 & & pollinates 15,854 \\
\midrule
Hierarchy & MONDO & GloBI taxonomy \\
\quad child--parent links & 17,801 & 77,852 \\
\quad families, $\ge$2 members & 3,182 & 8,370 \\
\quad median / max size & 3 / 150 & 3 / 315 \\
\midrule
Text & Wiki.\ + PubMed & Wikipedia \\
\quad sources with text & 574 drugs & 5,425 species \\
\quad facts measured & 3,444 & 23,694 \\
Questions & 856 & 900 \\
\bottomrule
\end{tabular}
\caption{Data for both domains. PrimeKG links drugs to diseases, and MONDO groups diseases into families. GloBI links species to species, and its taxonomy groups species into genera. A \emph{family} is a parent with two or more children. \emph{Facts measured} are edges we check for a mention of the answer.}
\label{tab:data}
\end{table}


\paragraph{Reference text corpora.}
For PrimeKG drugs, we use English Wikipedia articles and PubMed abstracts. For GloBI organisms, we use species' Wikipedia articles (dropping stubs $<$\num{3000} characters). All articles are stripped of markup, templates, infoboxes, tables, and file links.

\paragraph{Evidence granularity.}
We classify retrieved text into three levels: \emph{named} (exact name, synonym, abbreviation), \emph{family-only} (broad parent category within two levels), and \emph{silent} (neither). A control with unrelated disease pairs shows that matches track facts, not the size of the name set (unrelated matches occur in only 0.9--1.4\% of Wikipedia articles, versus 6.7--19.3\% for true diseases).

\paragraph{Corpus coverage.}
Evaluating 1,800 PrimeKG facts with Qwen2.5 0.5B (closed book) reveals that for 66.4\% of facts, the text lacks the exact disease name and the model lacks the fact. As detailed in Table~\ref{tab:boxes} and Table~\ref{tab:coverage}, the exact disease is missing in 84.9\% of Wikipedia articles, 81.1\% of PubMed sets, and 77.5\% combined, with random guessing accounting for most closed-book correct answers on unnamed facts.

\begin{table}[t]
\centering
\small
\setlength{\tabcolsep}{4pt}
\begin{tabular}{@{}lrrrrr@{}}
\toprule
 & & \multicolumn{2}{c}{Text names it} & \multicolumn{2}{c}{Text does not} \\
\cmidrule(lr){3-4}\cmidrule(l){5-6}
Relation & $n$ & knows & not & knows & not \\
\midrule
Indication & 600 & 9.7 & 7.2 & 21.8 & 61.3 \\
Contraind. & 600 & 1.5 & 5.2 & 20.7 & 72.7 \\
Off-label & 600 & 4.3 & 6.2 & 24.2 & 65.3 \\
\midrule
All & 1,800 & 5.2 & 6.2 & 22.2 & \textbf{66.4} \\
\bottomrule
\end{tabular}
\caption{Share of PrimeKG facts (\%) by whether the drug's Wikipedia article
names the disease and whether Qwen2.5 0.5B picks it from five options closed
book (``knows'').}
\label{tab:boxes}
\end{table}

\begin{table}[t]
\centering
\small
\setlength{\tabcolsep}{4pt}
\begin{tabular}{@{}llrrrr@{}}
\toprule
Relation & Text & $n$ & Named & Family & Silent \\
\midrule
\multirow{3}{*}{Indication} & Wikipedia & 1,722 & 20.3 & 15.5 & 64.2 \\
 & PubMed & 1,722 & 26.3 & 18.6 & 55.1 \\
 & both & 1,722 & 29.0 & 19.7 & 51.3 \\
\midrule
\multirow{3}{*}{Contraind.} & Wikipedia & 1,722 & 9.8 & 4.1 & 86.1 \\
 & PubMed & 1,722 & 11.4 & 5.4 & 83.2 \\
 & both & 1,722 & 15.9 & 6.9 & 77.2 \\
\midrule
PrimeKG & both & 3,444 & 22.5 & 13.3 & 64.2 \\
\midrule
Eats & Wikipedia & 8,000 & 4.0 & 6.6 & 89.3 \\
Preys on & Wikipedia & 8,000 & 2.2 & 4.4 & 93.4 \\
Parasite of & Wikipedia & 2,389 & 12.5 & 15.5 & 72.0 \\
Pollinates & Wikipedia & 5,305 & 2.2 & 8.0 & 89.7 \\
\midrule
GloBI & Wikipedia & 23,694 & 3.9 & 7.1 & 89.1 \\
\bottomrule
\end{tabular}
\caption{Share of knowledge-graph facts whose text names the answer, names only its family, or names neither. PrimeKG uses 574 drugs with three facts of each relation.}
\label{tab:coverage}
\end{table}


\paragraph{Closed-book recall.}
Although family-only text covers 13.3\% of PrimeKG facts, testing on Qwen2.5-7B shows that coarse family mentions do not reflect stronger underlying model knowledge: it answers 39.6\% of family-only facts, 32.0\% of silent facts and 65.7\% of named facts correctly (Appendix~B). Odds ratio (OR) compares the odds of a correct closed-book answer in two groups. OR$=$1 means no difference, and OR$<$1 means the first group is answered correctly less often. We pool the relations with the Mantel--Haenszel estimator; $p$ is the chance of a gap this large if the groups did not differ. Family-only facts are answered far less often than named facts ($\mathrm{OR}=0.27$, $p = 9\times10^{-9}$) and as often as silent facts ($\mathrm{OR}=1.10$, $p = 0.61$). While an exact Wikipedia mention yields a 3.7$\times$ increase in the odds of model knowledge, a family name provides no measurable advantage on its own. This motivates testing how models handle coarse partial evidence at inference time, and whether symbolic ontologies can recover exact answers.



